\chapter{複素引数凸関数}
    \section{複素引数凸関数の例}
        \label{複素引数凸関数の例}
        \subsection{\texorpdfstring{$\|\sum_{i=1}^N X_i\bm{a}_i + \bm{b}\|_2^2\;(\bm{a}_i, \bm{b} \in \complexNumbers^n,\;X_i \in \complexNumbers^{m\times n})$}{行列X1,X2,...,XNの線形結合×ベクトル+定数のノルム}は\texorpdfstring{$X_1,\dotsc,X_N$}{X1,X2,...,XN}に関して凸である。}
            \begin{proof}
                \quad\par
                $T_X \coloneqq (X_1,\dotsc,X_N),\;T_Y \coloneqq (Y_1,\dotsc,Y_N),\;\lambda \in [0,1]$とし、$f(T_X) \coloneqq \|\sum_{i=1}^N X_i\bm{a}_i + \bm{b}\|_2^2$とする。
                \begin{align*}
                    f((1-\lambda)T_X + \lambda T_Y) &= \left\|\sum_{i=1}^N ((1-\lambda)X_i + \lambda Y_i)\bm{a}_i + \bm{b}\right\|_2^2 \\
                    &= \left\|(1-\lambda)\left(\sum_{i=1}^N X_i\bm{a}_i + \bm{b}\right) + \lambda\left(\sum_{i=1}^N Y_i\bm{a}_i + \bm{b}\right)\right\|_2^2 \\
                    &= (1-\lambda)^2\left\|\sum_{i=1}^N X_i\bm{a}_i + \bm{b}\right\|_2^2 + \lambda^2\left\|\sum_{i=1}^N Y_i\bm{a}_i + \bm{b}\right\|_2^2 \\
                    &\phantom{=} + 2\lambda(1-\lambda)\Re{\left(\sum_{i=1}^N X_i\bm{a}_i + \bm{b}\right)^\mathrm{H}\left(\sum_{i=1}^N Y_i\bm{a}_i + \bm{b}\right)}
                \end{align*}
                よって
                \begin{align*}
                    &\phantom{=} (1-\lambda)f(T_X) + \lambda f(T_Y) - f((1-\lambda)T_X + \lambda T_Y) \\
                    &= \lambda(1-\lambda)\left[\left\|\sum_{i=1}^N X_i\bm{a}_i + \bm{b}\right\|_2^2 + \left\|\sum_{i=1}^N Y_i\bm{a}_i + \bm{b}\right\|_2^2 - 2\Re{\left(\sum_{i=1}^N X_i\bm{a}_i + \bm{b}\right)^\mathrm{H}\left(\sum_{i=1}^N Y_i\bm{a}_i + \bm{b}\right)}\right] \\
                    &= \lambda(1-\lambda)\left\|\left(\sum_{i=1}^N X_i\bm{a}_i + \bm{b}\right) - \left(\sum_{i=1}^N Y_i\bm{a}_i + \bm{b}\right)\right\|_2^2 \geq 0
                \end{align*}
            \end{proof}
        \subsection{\texorpdfstring{$\norm{A\bm{x}+\bm{b}}_2^2$}{||Ax+b||\string^2}が狭義凸\texorpdfstring{$\iff A^*A\succ O$}{⇔ A*Aが正定}}
            \begin{shadebox}
                $A \in \complexNumbers^{m\times n}, \bm{x} \in \complexNumbers^n, \bm{b} \in \complexNumbers^m$とする。
                $f(\bm{x}) \coloneqq \norm{A\bm{x}+\bm{b}}_2^2$が$\bm{x}$に関して狭義凸となる必要十分条件は$A^*A\succ O$である。
            \end{shadebox}
            $\lambda \in (0,1),\; \bm{x},\bm{y} \in \complexNumbers^m, \bm{x}\neq\bm{y}$とする。
            \begin{proof}
                まず次式が成り立つ。
                \begin{align*}
                    f((1-\lambda)\bm{x} + \lambda\bm{y}) &\coloneqq \norm{(1-\lambda)(A\bm{x}+\bm{b}) + \lambda(A\bm{y}+\bm{b})}_2^2 \\
                    &= (1-\lambda)^2 f(\bm{x}) + \lambda^2 f(\bm{y}) + 2\lambda(1-\lambda)\Re{(A\bm{x}+\bm{b})^*(A\bm{y}+\bm{b})}
                \end{align*}
                これより
                \begin{align*}
                    &\phantom{=} (1-\lambda)f(\bm{x}) + \lambda f(\bm{y}) - f((1-\lambda)\bm{x} + \lambda\bm{y}) \\
                    &= \lambda(1-\lambda)\left[f(\bm{x}) + f(\bm{y}) - 2\Re{(A\bm{x}+\bm{b})^*(A\bm{y}+\bm{b})}\right] \\
                    &= \lambda(1-\lambda)\norm{(A\bm{x}+\bm{b}) - (A\bm{y}+\bm{b})}_2^2 = \lambda(1-\lambda)\norm{A(\bm{x}-\bm{y})}_2^2 \\
                    &= \lambda(1-\lambda)(\bm{x}-\bm{y})^*A^*A(\bm{x}-\bm{y})
                \end{align*}
                $\lambda \in (0,1),\;\bm{x}\neq\bm{y}$に対して上式が正となる必要十分条件は$A^*A\succ O$である。
            \end{proof}